\documentclass{article}

% a bunch of useful packages.
\usepackage[inference]{semantic}
\usepackage{amsmath} \allowdisplaybreaks % lets align equations break over pages.
\usepackage{mathtools}
\usepackage{hyperref}
\usepackage[amsmath,hyperref,amsthm]{ntheorem}
\usepackage{thmtools}
\usepackage[capitalize]{cleveref}
\usepackage{mathpartir}
\usepackage{stmaryrd}
\usepackage{amssymb}
\usepackage{latexsym}
\usepackage{color}
\usepackage[dvipsnames]{xcolor} % the best way to colour text
\usepackage[colorinlistoftodos]{todonotes} 
\usepackage{tikz} \usetikzlibrary{positioning,shadows,arrows,calc,backgrounds,fit,shapes,shapes.multipart,decorations.pathreplacing,shapes.misc,patterns}
\usepackage{xspace}
\usepackage{scalerel}
%\usepackage{bm} %bold math. a mess to use.
\usepackage{bussproofs} % for logic-style proofs
\usepackage[
    left=1cm, 
    right=5cm,    % Wide margin for todo notes
    top=2cm, 
    bottom=2cm, 
    marginparwidth=4cm, % Must be smaller than the 'right' margin
    marginparsep=5mm    % Gap between main text and the note
]{geometry}
\hypersetup{ pdfpagemode=UseOutlines, colorlinks=true, linkcolor=green, citecolor=blue }

% organise your code.
\input{cmds}

% \newcommand{\box}[1]{\colorbox{green!20}{#1}}
% Elements of standard Wasm
\newcommand{\ws}[1]{\mathtt{#1}}
% Elements introduced by CWasm
\newcommand{\newt}[1]{\mathtt{\color{magenta}{#1}}}
\newcommand{\new}[1]{\textcolor{magenta}{#1}}

% Things I either still need to define or find a better symbol for
\newcommand{\colt}{\mathbb{N}}  % Color type

\newcommand{\numt}{\tau_n}      % Numeric type
\newcommand{\tagt}{\tau_{tag}}  % Tag type
\newcommand{\ptrt}{\tau_{ptr}}  % Pointer type
\newcommand{\addrt}{i32}

\newcommand{\tagnum}{\bigcirc}
\newcommand{\tagptr}{\square}

\newcommand{\te}[1]{\text{#1}}

\newcommand{\ang}[1]{\langle #1\rangle}
\newcommand{\thook}[1]{\xhookrightarrow{#1}}
\newcommand{\thookt}[1]{\thook{\text{#1}}}

\newcommand{\match}[5]{\mt{match}~#1\left|\begin{array}{ll}#2 &\Rightarrow #3 \\ #4 &\Rightarrow #5 \end{array}\right.}
\newcommand{\none}{\mt{None}}
\newcommand{\some}[1]{\mt{Some(#1)}}

\title{
    CWasm\\
    \large{(Diego's version)}
}
\author{Diego Oniarti}
\date{}

\begin{document}
\maketitle

\section{Values}
A new value is introduced that is the colored pointer (here referred to just pointer
for brevity). Pointers are kin to CHERI capabilities and have three components:
\begin{enumerate}
    \item addr: An address in memory
    \item color: Pointers are used to access the new ``colored memory". The color
        of a pointer must match the color of the memory region it tries to
        access, lest it traps.
    \item id An identifier given to the pointer upon allocation. This id is preserved
        when the pointer is incremented and decremented.
\end{enumerate}
For now we assume infinite possible colors, represented by natural numbers. In later
stages of the project I'm going to limit the number of colors.\\
More in detail, colors are \textit{options} of $\mathbb N$, where the \mt{None} case
is used to represent memory that hasn't been allocated yet or the color of pointers
generated by a failed allocation.

\begin{align*}
    && \ws{byte} \bnfdef&\ \ws{0x00} \mid \dots \mid \ws{0xff}
    \\
    \te{(integers)} && \left\{ \begin{aligned}
            \ws{uN} \bnfdef& \\
            \ws{sN} \bnfdef& \\
            \ws{iN} \bnfdef&
    \end{aligned} \right. &\
    \begin{aligned}
        & 0 \mid 1 \mid \dots \mid 2^N-1 \\
        & -2^{N-1} \mid \dots \mid -1 \mid 0 \mid 1 \mid \dots \mid 2^{N-1}-1 \\
        & \ws{uN}
    \end{aligned}
    \\
    && \ws{name} \bnfdef&\ \ws{char}^*
    \\
    && \ws{char} \bnfdef&\ \ws{U+00} \mid \dots \mid \ws{U+D7FF} \mid \ws{U+E000} \mid \dots \mid \ws{U+10FFFF}
    \\
    \te{(pointer)}&& \newt{p} \bnfdef&\ \newt{ \{addr: \addrt,\ color: \colt^?,\ id: \mathbb N\} }
    \\
    \te{(value)} && \ws{v} \bnfdef&\ \ws{\new{\numt}.const\ c} \mid \newt{\ptrt.const\ p}
    \\
                 &&& \dots \te{floats and others unchanged} \dots
\end{align*}

\section{Types}
A new type for pointers has been introduced ($\ptrt$), while the old numeric type get
relegated to the ``numeric type" $\numt$. The more generic $\tau$ takes both
classes of types.

\begin{align*}
    \te{(numeric)}
    &&\newt{\numt} \bnfdef&\
    \ws{i32} \mid \ws{i64} \mid \ws{f32} \mid \ws{f64}
    \\
    \te{(value type)}
    && \newt{\tau} \bnfdef&\
    \newt{\numt} \mid \newt{\ptrt}
    \\
    \te{(function/instructions)}
    &&\ws{\rho} \bnfdef&\
    \ws{\tau^* \to \tau^*}
    \\
    &&& \dots rest \dots
\end{align*}

\section{Runtime instances}
The ``colored memory" is introduced as an alternative to the standard heap. Bytes
in colored memory are ``tagged" with a color and a tag. The color is set to an area
of memory when it is allocated, and it is used to ensure spatial safety. The tag
distinguishes between numeric values and handles. New instructions are used to load
and store handles to memory, and a load will trap if the memory is not tag as
containing a pointer.

\begin{align*}
    \te{(heap)} && \ws{H} \bnfdef&\ \ws{\{data: byte^*,\ max: u32^?\}}
    \\
    \te{(tagged byte)} && \newt{tbyte} \bnfdef&\ \newt{\{val:byte, col: \colt^?, tag: \tagt\}}
    \\
    \textstyle \left( \begin{matrix} \text{colored memory} \\ \text{instance}\end{matrix} \right) && \newt{\Xi} \bnfdef&\ \newt{\{ data: tbyte^*,\ max : u32^?\}}
    \\
    \te{(allocator instance)} && \new{A} \bnfdef&\ \new{\emptyset} \mid \new{A} \newt{[id\mapsto c]}  % \new{id \to \colt^?}
    \\
    \te{(num/ptr)}
    && \newt{\tagt} \bnfdef&\
    \newt{\tagnum} \mid \newt{\tagptr}
    \\
    \te{(store)} && \ws{S} \bnfdef&\ \ws{\{funcs, tables, mems, globals,} \\
                 &&               & \newt{colored\_mem: \Xi}, \newt{allocator:A}\}
\end{align*}
The store gets extended with an instance of the colored memory and an instance of
the allocator.

The allocator is a map between all valid pointer ids and their color. When a pointer
gets freed, the entry in the allocator with its pointer is removed.

\section{Instructions}
Instructions defined with reference both to \href{https://www.w3.org/2021/11/wasm-stage/syntax/instructions.html}{Wasm} and to the Iris-MSWasm paper\cite{iriswasm}.

\subsection{Numeric instructions}
Numeric instructions are modified to only work on numeric types ($\numt$) instead
of all types ($\tau$).

Some new $relop$ instructions may be added for $\ptrt$. These would only compare the
address of the two pointers.

\subsection{Parametric}
Parametric instructions are unchanged.

\subsection{Variable instructions}
Variable instructions are also unchanged. Pointers can be set and gotten like any
other value.

\subsection{Administrative instructions}
The modifications to the administrative instructions are the same as in the\\ iris-MSWasm paper\cite{iriswasm}.
\begin{align*}
    \ws{ai} \bnfdef&\ \ws{bi} \mid \newt{\ptrt.const\ p} \mid \ws{trap} \mid \ws{invoke\ i} \mid \dots
\end{align*}
$\newt{\ptrt.const~p}$ is not allowed as a basic instruction since it would be a
trivial case of pointer forgery.

\subsection{Memory instructions}
The existing memory instructions work as usual for the types already existing in
Wasm. The $\ws{load}$ operation is not defined for the $\ptrt$ type, to avoid forgery.
\todo{$\newt{\ptrt.store}$ could still be allowed, but should it?}

\begin{align*}
    \te{(basic instructions)} && \ws{bi} \bnfdef&\ \dots \mid
    \newt{colalloc} \mid
    \newt{colfree} \mid
    \newt{\ptrt.add} \\ &&&
    \mid \newt{\tau.colstore} \mid
    \newt{\tau.colload}
\end{align*}

\begin{itemize}
    \item $\newt{colalloc}$: pops an int $n$ from the stack and allocate $n$ bytes
        of memory in the colored memory $\Xi$. The allocated bytes need to be zeroed
        and colored with a fresh color. The instruction puts a pointer $\newt{p}$ on
        the stack. The pointer has a fresh id and the same color as the memory.
    \item $\newt{colfree}$: pops a $\newt{\ptrt}$ from the stack. The id of the pointer
        is invalidated by the allocator, which will make all subsequent uses of it
        trap.
        \footnote{The instructions as described here do not take into consideration
            the allocator strategy. New allocations should be able to overwrite freed
        areas of memory}
    \item $\newt{\ptrt.add}$: pops a $p: \newt{\ptrt}$ and a $n: \ws{\addrt}$ from the
        stack. Adds $n$ to the address component of $p$. Incrementing (or
        decrementing) a pointer address outside of its allocated area does not cause
        a trap. Only the utilization of the pointer (load/store) will trap.
    \item $\newt{\tau.colstore}$: pops a pointer and a value from the stack. Interprets
        the value as $\tau$ and writes its byte representation to the colored memory.\\
        If $\tau$ is $\ptrt$, the value must be a $\newt{\ptrt.const\ p}$. If $\tau$
        is $\numt$, the value must be $\ws{\numt}$.\\
        All bytes being accessed in memory must have the same color as the pointer.
        The appropriate tag will be set the altered bytes, depending on $\tau$.
    \item $\newt{\tau.colload}$: Pops a pointer from the stack, loads the value in
        colored memory that pointer points to, and pushes the value onto the stack.
        If $\tau$ is $\ptrt$ all bytes read from memory need to have the $\tagptr$ tag.
\end{itemize}

\section{Instruction typing}
\begin{center}
    \typerule{wt-colalloc}{}{
        \Gamma \vdash \newt{colalloc}: [\addrt] \to [\newt{\ptrt}]
    }{colalloc}
    \typerule{wt-colfree}{}{
        \Gamma \vdash \newt{colfree}: [\newt{\ptrt}] \to []
    }{colfree}

    \typerule{wt-ptradd}{}{
        \Gamma \vdash \newt{\ptrt.add}: [\newt{\ptrt}, \addrt] \to [\newt{\ptrt}]
    }{ptradd}

    \typerule{wt-colstore}{}{
        \Gamma \vdash \newt{\numt.colstore}: [\newt{\ptrt}, \numt] \to []
    }{colstore}
    \typerule{wt-colstore-ptr}{}{
        \Gamma \vdash \newt{\ptrt.colstore}: [\newt{\ptrt}, \ptrt] \to []
    }{colstore-ptr}

    \typerule{wt-load}{}{
        \Gamma \vdash \newt{\numt.colload} [\newt{\ptrt}, \numt] \to [\numt]
    }{colload}
    \typerule{wt-load-ptr}{}{
        \Gamma \vdash \newt{\ptrt.colload} [\newt{\ptrt}, \ptrt] \to [\ptrt]
    }{colload-ptr}
\end{center}

\section{Allocator semantics}
Here are some helpers stating how \textit{allocation} and \textit{free} operations
alter colored memory and the allocator itself.
\begin{align*}
    % \new{valid\_ptr} \triangleq&~ A\mt{(p.id)} = \mt{Some(\_)}
    % \\
    % \new{color\_match(p, tbytes)} \triangleq&~ \forall \mt{tb} \in \mt{tbytes}, \mt{p.color}=\mt{tb.color}
    % \\
    %\newt{isfree(a)} \triangleq&~\Xi[\mt{a}]=\mt{None} \vee \nexists \mt{id}.A\mt{(id)}=\mt{Some(\_)}
    \newt{isfreshcol(c)} \triangleq&~ \nexists\mt{id}.A(\mt{id})=\mt{c}
    \\
    \newt{isfree(a)} \triangleq&~ \match{\Xi\mt{[a]}}
    {\mt{None}}{true}
    {\mt{Some(c)}}{\mt{isfreshcol(c)}}
    \\
    \newt{isfreshid(id)} \triangleq&~ A\mt{(id)}=\mt{None}
\end{align*}

\begin{center}
    \typerule{alloc}{
        \forall \mt{a\_off} \in [\mt{a},\mt{a+n}).\mt{isfree(a\_off)} &
        \mt{isfreshcol(c)} & \mt{isfreshid(id)} \\
        \Xi' = \Xi\{\mt{with data}=\Xi\mt{.data}[\mt{addr..addr+n}:=\ang{\mt{0x00},\mt{c},\tagnum}]\}\\
        A'=A'[\mt{id}\mapsto\mt{c}]
    }{
        \ang{\Xi,A}\thook{\newt{alloc(a,n,c,id)}}\ang{\Xi',A'}
    }{allocsem_alloc}

    \typerule{free}{
        A' = A[\mt{id}\mapsto\mt{None}]
    }{
        \ang{\Xi,A}\thook{\newt{free(id)}}\ang{\Xi,A'}
    }{allocasem_free}
\end{center}

Intuitively, allocation requires an area of memory to be free in the first place, and
for the new \mt{color} and \mt{id} to be fresh. After the allocation the area of memory
$[\mt{addr..addr+n})$ is zeroed, colored with the new color, and tagged as numeric.

\section{Operational semantics}
The judgement for the operational semantics of instructions assumes the form
$ (S, F, []) \hookrightarrow (S', F', []) $. This means a store, a frame, and a stack
step to a new configuration where one or more of the components is altered.
The stack contains administrative instructions and values (represented by the
$\tau.const~v$ administrative instruction themselves).

For ease of reading the store $S$ is exploded in its relevant components, transforming
the judgement in $ (\ang{\Xi, A}, F, []) \hookrightarrow (\ang{\Xi', A'}, F, []) $.

Some other notation abuse:
\begin{itemize}
    \item $\Xi[n] \triangleq \Xi\mt{.data}[n]$
    \item $|\Xi| \triangleq \match{\Xi.\mt{max}}{\none}{\infty}{\some{$n$}}{n}$
    \item $|\tau|$ is the number of bytes required to serialize values of type $\tau$
\end{itemize}

More helpers
\begin{align*}
    \newt{aligned(addr,\tau)} \triangleq&~ \exists k:\mathbb N.~\mt{addr} = |\tau|\cdot k \\
    \newt{serialize(val)} \triangleq&~ \mt{the byte representation of val} \\
    \newt{deserialize(bs,\tau)} \triangleq&~ \mt{the $\tau$ interpretation of the bytes bs}
\end{align*}

\begin{center}
    \typerule{opsem-colalloc}{
        \ang{\Xi,A}\thook{\mt{alloc(a,n,c,id)}}\ang{\Xi',A'} &
        \mt{p} = \{\mt{addr}=\mt{a}, \mt{color}=\mt{c}, \mt{id}=\mt{id}\}
    }{
        (\ang{\Xi, A}, F, [\ws{\addrt.const~n}, \newt{colalloc}]) \hookrightarrow (\ang{\Xi', A'}, F, [\newt{\ptrt.const~p}])
    }{opsem-colalloc}

    \typerule{opsem-colalloc-fail}{
        \mt{p} = \{\mt{addr}=0, \mt{color}=\mt{None}, \mt{id}=0\}
    }{
        (S, F, [\ws{\addrt.const~n}, \newt{colalloc}]) \hookrightarrow (S, F, [\newt{\ptrt.const~p}])
    }{opsem-colalloc-fail}

    \typerule{opsem-colfree}{
        \mt{p}=\{\mt{addr=a}, \mt{color=Some(c)}, \mt{id=n}\} &
        \mt{A(n)=c} &
        \ang{\Xi,A}\thook{\newt{free(n)}}\ang{\Xi,A'}
    }{
        (\ang{\Xi, A}, F, [\ptrt\mt{.const p}, \mt{colfree}]) \hookrightarrow (\ang{\Xi, A'}, F, [])
    }{opsem-colfree}

    \typerule{opsem-colstore}{
        % pointer parts
        \mt{p}=\{\mt{addr=a}, \mt{color=Some(c)}, \mt{id=n}\} &
        % Pointer hasn't been freed
        A\mt{(n)=c} \\
        % byte representation of the value
        \mt{serialise($\numt$, v)=bs} &
        % check bounds
        \mt{a+|bs|<$|\Xi|$} \\
        % color agreement
        \forall \mt{aoff} \in [\mt{a,a+|bs|}). \Xi\mt{[aoff].color=Some(c)} \\
        % copy the value
        \forall \mt{i} \in [\mt{0,|bs|}).~
        \Xi'[\mt{a+i}]=\{\mt{val=bs[i],col=Some(c),tag=}\tagnum\} % The color entry in the records is not updated. Maybe there is a way to make it more explicit
        % Memory bound
    }{
        (\ang{\Xi, A}, F, [\ptrt\mt{.const~p}, \numt\mt{.const~v}, \numt\mt{.colstore}]) \hookrightarrow (\ang{\Xi', A}, F, [])
    }{opsem-colstore}

    \typerule{opsem-colstore-ptr}{
        \mt{p}=\{\mt{addr=a}, \mt{color=Some(c)}, \mt{id=n}\} &
        A\mt{(n)=c} \\
        \mt{serialise($\ptrt$, v)=bs} &
        \mt{a+|bs|<$|\Xi|$} \\
        \forall \mt{aoff} \in [\mt{a,a+|bs|}). \Xi\mt{[aoff].color=Some(c)} \\
        \forall \mt{i} \in [\mt{0,|bs|}).~
        \Xi'[\mt{a+i}]=\{\mt{val=bs[i],col=Some(c),tag=}\tagptr\} \\
        \mt{aligned(a,$\ptrt$)}
    }{
        (\ang{\Xi, A},F,[\ptrt\mt{.const~p}, \ptrt\mt{.const~v}, \ptrt\mt{.colstore}]) \hookrightarrow (\ang{\Xi', A}, F, [])
    }{opsem-colstore-ptr}

    \typerule{opsem-colload}{
        % destruct the pointer fields
        \mt{p}=\{\mt{addr=a}, \mt{color=Some(c)}, \mt{id=n}\} &
        % check the pointer is valid
        A\mt{(n)=c} &
        % bounds check
        \mt{a+|$\tau$|<$|\Xi|$} \\
        % select the tagged bytes from colored memory
        \mt{tb}s = \Xi[\mt{a..a+|$\tau$|}] &
        % check that all have the right color and tag
        \forall \mt{tb} \in \mt{tb}s.~ \mt{tb.col}=\some{c} \wedge \mt{tb.tag}=\tagnum \\
        % strip the "extra" data from the tabbed bytes to get just the bytes
        \mt{b}s = [\mt{tb.val} \forall \mt{tb} \in \mt{tb}s] &
        % interpret the value
        \mt{deserialize(bs, $\numt$)} = \mt{v}
    }{
        (S,F,[\ptrt\mt{.const~p}, \numt\mt{.colload}]) \hookrightarrow (S, F, [\numt\mt{.const}~v])
    }{opsem-colload}

    \typerule{opsem-colload-ptr}{
        % destruct the pointer fields
        \mt{p}=\{\mt{addr=a}, \mt{color=Some(c)}, \mt{id=n}\} &
        % check the pointer is valid
        A\mt{(n)=c} &
        % bounds check
        \mt{a+|$\tau$|<$|\Xi|$} \\
        % select the tagged bytes from colored memory
        \mt{tb}s = \Xi[\mt{a..a+|$\tau$|}] &
        % check that all have the right color and tag
        \forall \mt{tb} \in \mt{tb}s.~ \mt{tb.col}=\some{c} \wedge \mt{tb.tag}=\tagptr \\
        % strip the "extra" data from the tabbed bytes to get just the bytes
        \mt{b}s = [\mt{tb.val} \forall \mt{tb} \in \mt{tb}s] &
        % interpret the value
        \mt{deserialize(bs, $\ptrt$)} = \mt{v} \\
        \mt{aligned(a,$\ptrt$)}
    }{
        (S,F,[\ptrt\mt{.const~p}, \ptrt\mt{.colload}]) \hookrightarrow (S, F, [\ptrt\mt{.const}~v])
    }{opsem-colload-ptr}

    \typerule{opsem-ptradd}{
        \mt{p}=\{\mt{addr=a}, \mt{color=Some(c)}, \mt{id=n}\} \\
        a' = a+v & 0\leq a' \\ % < |\Xi| \\
        \mt{p'}=\{\mt{p with addr$\coloneqq$a'}\} \\
    }{
        (S,F,[\ptrt\mt{.const~p}, \addrt\mt{.const}~v]) \hookrightarrow (S, F, [\ptrt\mt{.const~p'}])
    }{opsem-ptradd}
\end{center}

\bibliographystyle{alpha}
\bibliography{refs}
\end{document}
